\documentclass[11pt]{article}
\usepackage[bottom]{footmisc} 
\renewcommand{\thempfootnote}{\arabic{mpfootnote}}

\usepackage[hidelinks]{hyperref}

\usepackage[utf8]{inputenc}

\usepackage[T1]{fontenc}
\usepackage[ttdefault=true]{AnonymousPro}

\usepackage{multicol}
\setlength{\columnsep}{1cm}

\usepackage{abc}
\usepackage{graphicx}
\usepackage{musicography}

\usepackage[margin=0.5in]{geometry} 
\geometry{letterpaper}


\usepackage{fancyhdr} 
\pagestyle{fancy} % options: empty , plain , fancy
\renewcommand{\headrulewidth}{0pt} % customise the layout...
\lhead{}\chead{}\rhead{}
\lfoot{}\cfoot{\thepage}\rfoot{}


\begin{document}

\thispagestyle{empty}

\begin{abc}[name=kettle,options={-O= +c -F single.fmt}]
X:1
T:Molly Put the Kettle On
T:Polly Put the Kettle On
M:C|
K:D
a2 ab a2 ff | gfed cA A2 | a2 ab a2 gg | fded [df]2 fg | 
a^gab a=gfa | gfed cA A2 | a2 ab a2 gg | fded f2 ef ||
efed fdgg | fded cB A(3A/B/c/ | d2 ed fdgd | fded [df]2 (3ABc | 
dded fdgd | fded cB A(3A/B/c/ | d2 edfdgd | fded [df]4 |]
\end{abc}

%\vskip 0.5in
\noindent\makebox[\textwidth][c]{%
  \begin{minipage}{7.0in}
    \setlength{\parindent}{1.5em}


Polly Put the Kettle On (Molly Put the Kettle On) is a Mother Goose
children's rhyme. The `modern' Mother Goose dates back to the early
18th century when Charles Perrault's fairy tale collection, \textbf{Contes de
ma Mère l'Oye} was translated into English as \textbf{Tales of My Mother
Goose}. Like the Germanic Brothers Grimm stories, Mother Goose tales
were taken from oral tradition and traces of the stories can be found
in early 17th century French literature. Mother Goose tales are also
loosely linked to such historical figures as the wife of King Robert
II of France, known as \textit{Berthe pied d'oie} (Goose-Footed Bertha),
Charlemagne's mother Bertrade of Laon, \textit{regina pede aucae} (the
goose-foot queen), and even the biblical Queen of Sheba.\footnote{\url{https://en.wikipedia.org/wiki/Mother_Goose}}

The tune, Polly Put the Kettle On, appears in several early 19th
century British musicians' personal collections, including those of
T.J. Dixon (Lincolnshire, 1798), Rev. R. Harrison (Cumbria, c. 1815),
and J. Lishman (Lake District, c. 1825). In America, the tune is found
in the personal collections of Gurdon Trumbull (Stonington, CT, 1801)
and Ann Winnington (New York,
1810).\footnote{\url{https://tunearch.org/wiki/Annotation:Polly\_Put\_the\_Kettle\_On\_(1)}}

The nursery rhyme is mentioned in Charles Dickens' Barnaby Rudge
(1841) which takes place during the Gordon Riots of 1780.\footnote{The
  Gordon Riots were riots in London, June 2--9, 1780, in response to
  the passing of the Papist Act of 1778. This act was intended to
  officially reduce discrimination against British Catholics in the
  British Miliary thereby assisting in replenishing its numbers during
  the American War of Independence.}

The children's nursery rhyme most likely heard
today:\footnote{As sung by my mother.}
    {\parindent0pt
      \raggedright
      ~\\
Polly put the kettle on, \\
Polly put the kettle on, \\
Polly put the kettle on, \\
We'll all have tea. \\
~\\
Sukey take it off again, \\
Sukey take it off again, \\
Sukey take it off again, \\
They've all gone away.\footnote{Opie, Iona and Opie, Peter, \textbf{The Oxford Dictionary of Nursery
  Rhymes} (Oxford University Press, 1951, 2nd ed., 1997),
  pp. 353–54.} \footnote{In mid-eighteenth century British
  middle-class families, `Polly' was a nickname for Mary and `Sukey'
  was a nickname for
  Susan. (\url{https://en.wikipedia.org/wiki/Polly\_Put\_the\_Kettle\_On})}
  ~\\
    }

    ---Kitty Steetle
    
  \end{minipage}
  \vspace*{\fill}
}

\end{document}
